\documentclass[final]{beamer}

\usepackage[size=a1,scale=1.]{beamerposter}

\usetheme{default} % or metropolis-like styling later

% Colors
\definecolor{MyBlue}{RGB}{25,70,140}


\setbeamercolor{headline}{bg=MyBlue, fg=white}

\setbeamertemplate{headline}{
  \begin{beamercolorbox}[wd=\paperwidth,ht=6cm,center]{headline}
    \vspace{0cm}
    {\usebeamerfont{title}\color{white}\inserttitle\par}
    \vspace{0.5cm}
    {\usebeamerfont{author}\color{white}\insertauthor\par}
    {\usebeamerfont{institute}\color{white}\insertinstitute\par}
    \vspace{1cm}
  \end{beamercolorbox}
}

\setbeamertemplate{block begin}{
  \vspace{1.5ex}
  \noindent
  {\usebeamerfont{block title}\bfseries
   \usebeamercolor[fg]{block title}\insertblocktitle}
  \par
  % \vspace{0.1ex}
  \noindent\rule{\linewidth}{2.pt}
  % \vspace{0.8ex}
  \usebeamerfont{block body}
}
\setbeamertemplate{block end}{
  \vspace{1.5ex}
}

\setlength{\parskip}{1em}
\setbeamerfont{block body}{size=\large}
\setbeamerfont{block title}{size=\huge}

\setbeamercolor{block alerted title}{fg=black}
\setbeamercolor{block alerted body}{bg=gray!15, fg=black}

\setbeamertemplate{block alerted begin}{
  \vspace{1ex}
  \begin{beamercolorbox}[wd=\linewidth,sep=1em]{block alerted body}
  {\bfseries\insertblocktitle\par}
  \vspace{0.5ex}
}

\setbeamertemplate{block alerted end}{
  \end{beamercolorbox}
  \vspace{1ex}
}

\setlength{\columnsep}{2cm}

\setbeamerfont{title}{size=\Huge,series=\bfseries}
\setbeamerfont{author}{size=\Large}
\setbeamerfont{institute}{size=\large}

% Title info
\title{Reward Misspecification and Its Impact on Learning Dynamics}

\author{Victor Schmit}
\institute{Centrale Supélec}

\begin{document}

\begin{frame}[t]

	% \maketitle
	\begin{columns}[t]

		% -------- Column 1 --------
		\begin{column}{0.32\textwidth}
			\begin{block}{Introduction}
				\begin{itemize}
					\item Reinforcement learning (RL) agents are designed to optimize cumulative reward, yet their performance critically depends on how this reward is specified. In this project, we investigate the impact of reward misspecification on the behavior of tabular RL algorithms. Using a simple gridworld environment, we compare the learning dynamics and resulting policies of Q-learning and SARSA under both correctly specified and intentionally biased reward functions.

					      \medskip

					\item Our experiments show that even small modifications to the reward structure can lead to significantly different behaviors, including the emergence of suboptimal or unintended strategies. In particular, we observe that off-policy methods such as Q-learning tend to exploit reward inconsistencies more aggressively, while on-policy methods like SARSA exhibit more conservative behavior under uncertainty.

					      \medskip

					\item These results highlight a fundamental limitation of reinforcement learning: agents optimize the reward signal as given, rather than the true objective intended by the designer. This study emphasizes the importance of careful reward design and provides insight into how different algorithms respond to misspecified objectives.
				\end{itemize}
			\end{block}

			\begin{block}{Methods}

				Your methods
			\end{block}
		\end{column}

		% -------- Column 2 --------
		\begin{column}{0.32\textwidth}
			\begin{block}{Results}

				Your results
			\end{block}
			\begin{alertblock}{Important Result}
				This stands out visually.
			\end{alertblock}
		\end{column}

		% -------- Column 3 --------
		\begin{column}{0.32\textwidth}
			\begin{block}{Conclusion}
				Your conclusion
			\end{block}
		\end{column}

	\end{columns}

\end{frame}

\end{document}
